\documentclass[UTF8]{ctexart}

\usepackage{amsmath, amssymb, amsthm}
\usepackage{geometry}
\geometry{a4paper, margin=2.5cm}
\usepackage{hyperref}

\newtheorem{theorem}{定理}[section]
\newtheorem{definition}[theorem]{定义}
\newtheorem{lemma}[theorem]{引理}

\title{数学笔记}
\author{Mifansans}
\date{\today}

\begin{document}
\maketitle
\tableofcontents

\section{微积分}

\subsection{高斯积分}

\begin{theorem}[高斯积分]
\[
\int_{-\infty}^{+\infty} e^{-x^2}\,dx = \sqrt{\pi}.
\]
\label{thm:gauss}
\end{theorem}

\begin{proof}
考虑二重积分
\[
\left(\int_{-\infty}^{+\infty} e^{-x^2}\,dx\right)^2
=
\iint_{\mathbb{R}^2} e^{-(x^2+y^2)}\,dx\,dy,
\]
再使用极坐标变换即可。
\end{proof}

由定理 \ref{thm:gauss} 可知该积分收敛。

\end{document}
